Soit $ABC$ un triangle équilatéral de côté $a$. On construit $\mathcal{C}_A$ le cercle contenu dans le plan médiateur de $[BC]$, de centre le milieu de $[BC]$, de rayon $r$ fixé. On construit de même $\mathcal{C}_B$ et $\mathcal{C}_C$ sur les côtés $[AC]$ et $[AB]$. Soit $A'$ un point de $\mathcal{C}_A$ tel qu'il existe deux points $B'$ et $C''$ de $\mathcal{C}_B$ à distance $a$ de $A'$. On note $C'$ le point de $\mathcal{C}_C$ symétrique de $C''$ par rapport au plan médiateur de $[BC]$. Montrer que $A'B'C'$ est un triangle équilatéral.
